\documentclass[aspectratio=169]{beamer}
\usepackage[utf8]{inputenc}
\usepackage[T1]{fontenc}
\usepackage{booktabs}
\IfFileExists{theme_finance_mba.sty}{\usepackage{theme_finance_mba}}{\usepackage{beamer/theme_finance_mba}}

\title[Investments Day 1]{Finance MBA -- Investments and Valuation: Day 1}
\author{Instructor Name}
\date{2027 Edition}

\begin{document}

\begin{frame}
  \titlepage
\end{frame}

\section{1}

\begin{frame}{Course Roadmap for Day 1}
\begin{itemize}
  \item Returns and risk: intuition and measurement
  \item Hands-on assignment with two-stock starter sheet (APPLE, AMAZON)
  \item Scaling to a broader market set
  \item Gold versus stocks and bonds: role in portfolios
\end{itemize}
\end{frame}

\begin{frame}{Assignment Submission Rules}
\begin{assignmentblock}{Formatting}
\begin{itemize}
  \item File name: \texttt{TeamName\_A1\_Day1.pdf}
  \item First line inside file: assignment number and day
  \item Include all team member names on the first slide/page
  \item Submit via the agreed LMS channel before class end
\end{itemize}
\end{assignmentblock}
\end{frame}

\begin{frame}{Transition: From Market Data to Portfolio Logic}
\begin{takeawayblock}
First, we compute clean return series on a simple data set. Then, we expand the same logic to the full universe of assets.
\end{takeawayblock}
\end{frame}

\begin{frame}{Assignment 1 Setup (Simplified Starter)}
\begin{itemize}
  \item Use only two assets first: APPLE and AMAZON
  \item Keep exactly 10 years of data in the worksheet
  \item Compute:
    \begin{itemize}
      \item simple and log returns,
      \item average returns,
      \item volatility,
      \item covariance and correlation.
    \end{itemize}
  \item Validate formulas before scaling up.
\end{itemize}
\end{frame}

\begin{frame}{Gold Versus Stocks and Bonds}
\begin{tabular}{lccc}
\toprule
Asset class & Return focus & Volatility & Inflation hedge \\
\midrule
Equities & Growth & High & Partial \\
Bonds & Income/stability & Low--medium & Weak--moderate \\
Gold & Diversification & Medium--high & Often stronger \\
\bottomrule
\end{tabular}

\vspace{0.75em}
\begin{takeawayblock}
Gold can improve diversification, but its standalone volatility and long flat periods require context.
\end{takeawayblock}
\end{frame}

\begin{frame}{End-of-Day Recap}
\begin{itemize}
  \item We built returns from raw prices in a transparent way.
  \item We practiced on a two-stock example before full-scale computation.
  \item We compared gold with equities and bonds in portfolio terms.
\end{itemize}
\end{frame}

\end{document}
